\chapter{Limitations}
\label{ch:limitations}

The answer in Chapter~\ref{ch:headline} is narrower than it sounds, and the
limits are part of the answer rather than a disclaimer appended to it. They are
in the same document as the result, and the most important ones are repeated in
the same paragraph as the result where they first appear.

\section{What the accuracy result does not license}

\textbf{It does not license any claim that language models are bad at this
task.} What was measured is one twelve billion parameter model at four bit
quantization, with one prompt, no fine tuning, on a synthetic corpus, where the
entire input is a short list of attribute strings.

That input shape is close to the worst case for a system whose advantage is world
knowledge, and close to the best case for one that memorises naming conventions.
A language model asked to reason about a page would have the page. This one had a
pruned descriptor, because giving it the raw document while the n-gram model got
a feature vector would have measured something else entirely.

\textbf{The two classical engines were fitted on this corpus's own training
split. The language model met the corpus for the first time at inference.} That
is a real asymmetry and it points in the direction of the observed result, not
against it. A fine tuned model, a better prompt, a larger model, a higher
precision quantization, or real pages could each move the result, and none of
them was varied.

\textbf{One prompt was tested.} There is no prompt sweep, and there deliberately
is not one, because a prompt sweep chosen after seeing the first result is a
search for a number. Testing several prompts is future work and it requires a
pre-registration of its own.

\section{The corpus is synthetic, and that is the load bearing limitation}

Behaviour on real world pages is unmeasured. Every number in this report is a
measurement on forms this repository generated.

The corpus was designed to be adversarial in the places that matter, and the
markup quality tiers exist precisely so that the easy case does not dominate the
average. But a generated corpus has a generator's habits, and a classifier fitted
on it learns those habits along with the phenomenon. The rule table has the same
exposure in a different form: its vocabulary was written by the same author as
the locale profiles that produce the label text it matches.

That shared authorship is the strongest single reason to distrust the accuracy
gap between the classical engines and the language model, and it is why the
runtime finding is presented as the more robust one.

\section{Six locales is not multilingual}

Six locales, one of them held out. Every locale in the corpus uses either Latin
or Japanese script. There is no right to left locale, no Indic script, no Cyrillic
locale, and no locale whose address grammar differs from the ones present in ways
the generator does not model.

The normalisation layer was hardened against the whole basic multilingual plane
after two property test failures outside the corpus's own alphabet, so the
tokeniser is not the weak point. The locale profiles are, and they are hand
written per locale, which means adding a locale is a real piece of work rather
than a configuration change.

\section{What the tool cannot see}

Repeated here from Chapter~\ref{ch:extractor} because a reader who skipped to the
results should still meet them:

\begin{itemize}
  \item Closed shadow roots are invisible to any script on the page, including
    this one.
  \item Cross origin frames are detected but not read. A hosted payment field is
    a common and correct pattern, and the finding says so rather than implying a
    defect.
  \item Canvas rendered form like widgets are reported as undetectable rather
    than guessed at.
  \item Only the step of a multi step form that is currently in the document is
    audited. There is no click through, because clicking is filling.
  \item \textbf{The tool does not verify that a browser actually autofills.} It
    verifies that the markup gives the browser what it needs. Those are different
    claims and only the second one is made.
\end{itemize}

\section{The evaluation's own limitations}

\textbf{The design resolves large effects, not small ones.} \vClusters{}
clusters give \vArrangements{} arrangements. That is enough to certify the
differences reported as certified in Chapter~\ref{ch:statistics} and not enough
to resolve the ones reported as not certified. More templates per family is the
only thing that helps; more resamples do not, because the test is already
exhaustive.

\textbf{Two statistical components run on this side of the package boundary.}
The clustered null and the absolute practical gate are computed here rather than
in the harness, because the harness has no public entry point for either. Both
are labelled in every result file they produce. The alternative, degrading
silently to unclustered resampling, would have made every comparison look more
significant than it is, and it was the one outcome ruled out before the phase
began.

\textbf{The evaluation runner classifies every form twice.} This is a real defect
and it is documented here rather than in an issue nobody reads. The runner
classifies once directly, which is where the run log's predicted label and per
field latency come from, and once through the audit engine, which is where the
finding codes come from. For the two deterministic engines the passes agree by
construction and cost microseconds. For the language model it doubled the
inference: the run made twice as many requests as there are forms, and about half
its wall clock is the second pass.

Two consequences follow, and the second matters to a reader of
Table~\ref{tab:findings}:

\begin{itemize}
  \item The language model benchmark cost about twice what it needed to.
  \item \textbf{A row's predicted label and its finding codes come from two
    different draws of a model that is not guaranteed to be deterministic.}
    Measured on this run, using the deterministic rule engine as a control for
    the check's own edge cases, the excess disagreement between the two passes is
    on the order of two fields in seven hundred. Small, real, and it should be
    known before anyone builds on the finding level numbers.
\end{itemize}

It is not fixed in this report's phase, and the reason is a rule rather than a
preference: fixing it changes what the evaluation runner does for every engine,
and doing that after seeing the results is the regeneration that this project's
reproducibility rules forbid. The fix is in Chapter~\ref{ch:future} and it is to
pass the first pass's predictions into the audit rather than letting the audit
reclassify, which also turns the two numbers into one fact.

\section{Development split evidence is one notch weaker}

Stated in full in Chapter~\ref{ch:method} and repeated here so that it is met
twice: the training manifest of the shipped model is marked as produced from a
working tree with uncommitted changes. The model artefacts and the corpus are
content addressed into every evaluation manifest, so the identity of what was
measured is not in doubt, but the development split tables in
Chapter~\ref{ch:classifiers} rest on a run whose tree state was not pinned to a
commit. They are reported and they are flagged.

\section{Class imbalance and small denominators}

Several labels in the taxonomy are carried by one template in one position, so
their per label rows in Table~\ref{tab:per-label} rest on small denominators.
Macro-F1 weights every label equally, which means those rows have the same weight
in the headline average as the labels that appear on every form.

That is a deliberate choice rather than an oversight: a label weighted average is
the one that notices when an engine quietly stops predicting a rare class. It
does mean the macro average is noisier than the micro average, and it is another
reason both are reported.

\section{Cost was not measured because nothing was billed}

The monetary cost column is null for every engine. The two classical engines have
no provider and the language model ran locally. Zero is a measurement and the
electricity was not free; null is the absence of one. No vendor price page was
retrieved during this work, so no list price estimate is computed anywhere, and
the price table document in this repository contains no prices on purpose. If a
cost column with a number in it is wanted, that is a retrieval with a date, not a
recollection.
